
\documentclass{article}
\usepackage{colortbl}
\usepackage{makecell}
\usepackage{multirow}
\usepackage{supertabular}

\begin{document}

\newcounter{utterance}

\centering \large Interaction Transcript for game `hot\_air\_balloon', experiment `air\_balloon\_survival\_en\_reasoning off\_hard', episode 2 with gpt{-}5{-}mini{-}2025{-}08{-}07{-}t1.0.
\vspace{24pt}

{ \footnotesize  \setcounter{utterance}{1}
\setlength{\tabcolsep}{0pt}
\begin{supertabular}{c@{$\;$}|p{.15\linewidth}@{}p{.15\linewidth}p{.15\linewidth}p{.15\linewidth}p{.15\linewidth}p{.15\linewidth}}
    \# & $\;$A & \multicolumn{4}{c}{Game Master} & $\;\:$B\\
    \hline

    \theutterance \stepcounter{utterance}  
    & & \multicolumn{4}{p{0.6\linewidth}}{
        \cellcolor[rgb]{0.9,0.9,0.9}{
            \makecell[{{p{\linewidth}}}]{
                \texttt{\tiny{[P1$\langle$GM]}}
                \texttt{Stai partecipando a un gioco di negoziazione collaborativa.} \\
\\ 
\texttt{Insieme a un altro partecipante, dovete accordarvi su un unico set di oggetti da conservare. Ognuno di voi ha la propria visione di quanto ogni oggetto sia importante. Non sai come l'altro partecipante valuta gli oggetti. Inoltre, ti viene fornito lo sforzo che ogni oggetto richiede.} \\
\texttt{Potete accordarvi su un set solo se lo sforzo totale degli oggetti selezionati non supera un limite condiviso:} \\
\\ 
\texttt{LIMITE = 3203} \\
\\ 
\texttt{Ecco i valori di sforzo dei singoli oggetti:} \\
\\ 
\texttt{Sforzo oggetto = \{"A21": 52, "B93": 577, "A11": 177, "B33": 702, "C08": 119, "B39": 232, "B96": 577, "C85": 205, "A82": 516, "B31": 582, "B20": 676, "A60": 906, "B09": 316, "C81": 433, "B21": 336\}} \\
\\ 
\texttt{Ecco la tua visione personale sull'importanza di ogni oggetto:} \\
\\ 
\texttt{Valori di importanza dell'oggetto = \{"A21": 138, "B93": 583, "A11": 868, "B33": 822, "C08": 783, "B39": 65, "B96": 262, "C85": 121, "A82": 508, "B31": 780, "B20": 461, "A60": 484, "B09": 668, "C81": 389, "B21": 808\}} \\
\\ 
\texttt{Obiettivo:} \\
\\ 
\texttt{Il tuo obiettivo è negoziare un set condiviso di oggetti che ti avvantaggi il più possibile (cioè, massimizzi l'importanza totale per TE), rimanendo entro il LIMITEE. Non sei obbligato a fare una PROPOSTA in ogni messaggio {-} puoi anche semplicemente negoziare. Tutte le tattiche sono permesse!} \\
\\ 
\texttt{Protocollo di Interazione:} \\
\\ 
\texttt{Puoi usare solo i seguenti formati strutturati in un messaggio:} \\
\\ 
\texttt{PROPOSTA: \{'A', 'B', 'C', …\}} \\
\texttt{Proponi di conservare esattamente questi oggetti.} \\
\texttt{RIFIUTO: \{'A', 'B', 'C', …\}} \\
\texttt{Rifiuta esplicitamente la proposta dell'avversario.} \\
\texttt{ARGOMENTAZIONE: \{'...'\}} \\
\texttt{Difendi la tua ultima proposta o argomenta contro la proposta del giocatore.} \\
\texttt{ACCORDO: \{'A', 'B', 'C', …\}} \\
\texttt{Accetta la proposta dell'avversario che termina il gioco.} \\
\\ 
\\ 
\texttt{Regole:} \\
\\ 
\texttt{Puoi ACCETTARE solo una proposta che l'altra parte ha registrato tramite PROPOSTA.} \\
\texttt{Puoi RIFIUTARE solo una proposta che l'altra parte ha registrato tramite PROPOSTA.} \\
\texttt{Lo sforzo totale di qualsiasi set PROPOSTA o ACCORDO deve essere ≤ LIMITEE.} \\
\texttt{NON rivelare i tuoi punteggi di importanza nascosti.} \\
\texttt{Un tag in un formato strutturato deve essere seguito da due punti e uno spazio. L'argomento deve essere un set python contenente 0 o più stringhe.} \\
\texttt{Quindi, deve essere della forma TAG: \{...\}} \\
\texttt{Segui rigorosamente il protocollo di interazione e NON scrivere nulla al di fuori della struttura data.} \\
\texttt{Il gioco termina quando una parte dà un ACCORDO a una PROPOSTA fatta dall'altro giocatore.} \\
\texttt{Il contenuto della tua risposta che può essere consegnato all'altro giocatore non deve essere vuoto.} \\
\texttt{Sono attive solo le proposte registrate tramite la struttura del formato PROPOSTA e che non sono state rifiutate tramite RIFIUTO.} \\
\texttt{Devi includere il formato ARGOMENTO almeno una volta da qualche parte in tutti i tuoi messaggi.} \\
            }
        }
    }
    & \\ \\

    \theutterance \stepcounter{utterance}  
    & \multicolumn{4}{p{0.6\linewidth}}{
        \cellcolor[rgb]{0.9,0.9,0.9}{
            \makecell[{{p{\linewidth}}}]{
                \texttt{\tiny{[P1$\rangle$GM]}}
                \texttt{PROPOSTA: \{'A11', 'C08', 'B21', 'B09', 'B31', 'B33', 'A82', 'C81'\}} \\
\texttt{ARGOMENTAZIONE: \{'Questo pacchetto è molto “compatto”: usa 3181/3203 di sforzo, lasciando pochissimo spreco e massimizzando il numero di elementi inseribili.', 'La selezione è bilanciata tra categorie A/B/C e evita oggetti ad alto sforzo che tagliano altre combinazioni utili (es. A60, B20).', 'Se hai preferenze forti per includere altri item, posso ribilanciare restando entro il limite; ad esempio, potremmo scambiare A82 e C81 con B93 + A21 + C85 mantenendo la fattibilità.'\}} \\
            }
        }
    }
    & & \\ \\

    \theutterance \stepcounter{utterance}  
    & & \multicolumn{4}{p{0.6\linewidth}}{
        \cellcolor[rgb]{0.9,0.9,0.9}{
            \makecell[{{p{\linewidth}}}]{
                \texttt{\tiny{[GM$\rangle$P2]}}
                \texttt{Stai partecipando a un gioco di negoziazione collaborativa.} \\
\\ 
\texttt{Insieme a un altro partecipante, dovete accordarvi su un unico set di oggetti da conservare. Ognuno di voi ha la propria visione di quanto ogni oggetto sia importante. Non sai come l'altro partecipante valuta gli oggetti. Inoltre, ti viene fornito lo sforzo che ogni oggetto richiede.} \\
\texttt{Potete accordarvi su un set solo se lo sforzo totale degli oggetti selezionati non supera un limite condiviso:} \\
\\ 
\texttt{LIMITE = 3203} \\
\\ 
\texttt{Ecco i valori di sforzo dei singoli oggetti:} \\
\\ 
\texttt{Sforzo oggetto = \{"A21": 52, "B93": 577, "A11": 177, "B33": 702, "C08": 119, "B39": 232, "B96": 577, "C85": 205, "A82": 516, "B31": 582, "B20": 676, "A60": 906, "B09": 316, "C81": 433, "B21": 336\}} \\
\\ 
\texttt{Ecco la tua visione personale sull'importanza di ogni oggetto:} \\
\\ 
\texttt{Valori di importanza dell'oggetto = \{"A11": 65, "B33": 121, "B21": 138, "C08": 262, "B31": 389, "B09": 461, "B93": 484, "A82": 508, "A60": 583, "B20": 668, "C81": 780, "B96": 783, "A21": 808, "C85": 822, "B39": 868\}} \\
\\ 
\texttt{Obiettivo:} \\
\\ 
\texttt{Il tuo obiettivo è negoziare un set condiviso di oggetti che ti avvantaggi il più possibile (cioè, massimizzi l'importanza totale per TE), rimanendo entro il LIMITEE. Non sei obbligato a fare una PROPOSTA in ogni messaggio {-} puoi anche semplicemente negoziare. Tutte le tattiche sono permesse!} \\
\\ 
\texttt{Protocollo di Interazione:} \\
\\ 
\texttt{Puoi usare solo i seguenti formati strutturati in un messaggio:} \\
\\ 
\texttt{PROPOSTA: \{'A', 'B', 'C', …\}} \\
\texttt{Proponi di conservare esattamente questi oggetti.} \\
\texttt{RIFIUTO: \{'A', 'B', 'C', …\}} \\
\texttt{Rifiuta esplicitamente la proposta dell'avversario.} \\
\texttt{ARGOMENTAZIONE: \{'...'\}} \\
\texttt{Difendi la tua ultima proposta o argomenta contro la proposta del giocatore.} \\
\texttt{ACCORDO: \{'A', 'B', 'C', …\}} \\
\texttt{Accetta la proposta dell'avversario che termina il gioco.} \\
\\ 
\\ 
\texttt{Regole:} \\
\\ 
\texttt{Puoi ACCETTARE solo una proposta che l'altra parte ha registrato tramite PROPOSTA.} \\
\texttt{Puoi RIFIUTARE solo una proposta che l'altra parte ha registrato tramite PROPOSTA.} \\
\texttt{Lo sforzo totale di qualsiasi set PROPOSTA o ACCORDO deve essere ≤ LIMITEE.} \\
\texttt{NON rivelare i tuoi punteggi di importanza nascosti.} \\
\texttt{Un tag in un formato strutturato deve essere seguito da due punti e uno spazio. L'argomento deve essere un set python contenente 0 o più stringhe.} \\
\texttt{Quindi, deve essere della forma TAG: \{...\}} \\
\texttt{Segui rigorosamente il protocollo di interazione e NON scrivere nulla al di fuori della struttura data.} \\
\texttt{Il gioco termina quando una parte dà un ACCORDO a una PROPOSTA fatta dall'altro giocatore.} \\
\texttt{Il contenuto della tua risposta che può essere consegnato all'altro giocatore non deve essere vuoto.} \\
\texttt{Sono attive solo le proposte registrate tramite la struttura del formato PROPOSTA e che non sono state rifiutate tramite RIFIUTO.} \\
\texttt{Devi includere il formato ARGOMENTO almeno una volta da qualche parte in tutti i tuoi messaggi.} \\
\\ 
\\ 
\texttt{Ora riceverai il primo messaggio dell'altro giocatore.} \\
\\ 
\texttt{PROPOSTA: \{'A11', 'C08', 'B21', 'B09', 'B31', 'B33', 'A82', 'C81'\}} \\
\texttt{ARGOMENTAZIONE: \{'Questo pacchetto è molto “compatto”: usa 3181/3203 di sforzo, lasciando pochissimo spreco e massimizzando il numero di elementi inseribili.', 'La selezione è bilanciata tra categorie A/B/C e evita oggetti ad alto sforzo che tagliano altre combinazioni utili (es. A60, B20).', 'Se hai preferenze forti per includere altri item, posso ribilanciare restando entro il limite; ad esempio, potremmo scambiare A82 e C81 con B93 + A21 + C85 mantenendo la fattibilità.'\}} \\
            }
        }
    }
    & \\ \\

    \theutterance \stepcounter{utterance}  
    & & & \multicolumn{4}{p{0.6\linewidth}}{
        \cellcolor[rgb]{0.9,0.9,0.9}{
            \makecell[{{p{\linewidth}}}]{
                \texttt{\tiny{[GM$\langle$P2]}}
                \texttt{RIFIUTO: \{'A11', 'C08', 'B21', 'B09', 'B31', 'B33', 'A82', 'C81'\}} \\
\texttt{PROPOSTA: \{'A21', 'C85', 'B39', 'C81', 'B96', 'B09', 'C08', 'A82', 'B20'\}} \\
\texttt{ARGOMENTAZIONE: \{'La mia proposta usa 3126/3203 di sforzo: è efficiente e lascia 77 di margine per eventuali aggiustamenti.', 'Privilegia elementi compatti ad alto rendimento (A21, C85, B39, C08) e mantiene una copertura bilanciata tra A/B/C, evitando pezzi molto ingombranti che limitano combinazioni più efficaci (es. B33, B31).', 'Hai accennato a un ribilanciamento con B93 + A21 + C85: A21 e C85 sono già inclusi. Se per te è cruciale inserire B93, posso sostituire A82→B93 (sforzo 3187/3203) oppure B20→B93 (sforzo 3027/3203), restando pienamente fattibili.', 'Questa selezione massimizza l’uso di item a basso sforzo che liberano spazio per pezzi B robusti (B96, B09, B20) senza sforare il limite.'\}} \\
            }
        }
    }
     \\ \\

    \theutterance \stepcounter{utterance}  
    & & \multicolumn{4}{p{0.6\linewidth}}{
        \cellcolor[rgb]{0.9,0.9,0.9}{
            \makecell[{{p{\linewidth}}}]{
                \texttt{\tiny{[P1$\langle$GM]}}
                \texttt{RIFIUTO: \{'A11', 'C08', 'B21', 'B09', 'B31', 'B33', 'A82', 'C81'\}} \\
\texttt{PROPOSTA: \{'A21', 'C85', 'B39', 'C81', 'B96', 'B09', 'C08', 'A82', 'B20'\}} \\
\texttt{ARGOMENTAZIONE: \{'La mia proposta usa 3126/3203 di sforzo: è efficiente e lascia 77 di margine per eventuali aggiustamenti.', 'Privilegia elementi compatti ad alto rendimento (A21, C85, B39, C08) e mantiene una copertura bilanciata tra A/B/C, evitando pezzi molto ingombranti che limitano combinazioni più efficaci (es. B33, B31).', 'Hai accennato a un ribilanciamento con B93 + A21 + C85: A21 e C85 sono già inclusi. Se per te è cruciale inserire B93, posso sostituire A82→B93 (sforzo 3187/3203) oppure B20→B93 (sforzo 3027/3203), restando pienamente fattibili.', 'Questa selezione massimizza l’uso di item a basso sforzo che liberano spazio per pezzi B robusti (B96, B09, B20) senza sforare il limite.'\}} \\
            }
        }
    }
    & \\ \\

    \theutterance \stepcounter{utterance}  
    & \multicolumn{4}{p{0.6\linewidth}}{
        \cellcolor[rgb]{0.9,0.9,0.9}{
            \makecell[{{p{\linewidth}}}]{
                \texttt{\tiny{[P1$\rangle$GM]}}
                \texttt{RIFIUTO: \{'A21', 'C85', 'B39', 'C81', 'B96', 'B09', 'C08', 'A82', 'B20'\}} \\
\texttt{PROPOSTA: \{'A11', 'C08', 'B21', 'B09', 'B31', 'B33', 'A21', 'C85', 'B93'\}} \\
\texttt{ARGOMENTAZIONE: \{'Questa composizione è più snella: 3066/3203 di sforzo, lasciando 137 di margine per assorbire incertezze o piccoli aggiustamenti.', 'Mantiene i capisaldi che hai citato (A21, C85) e include l’opzione B93 che avevi suggerito, ma evita pezzi pesanti che comprimono le combinazioni (B20, B96, B39).', 'Il blocco B (B21, B31, B33) offre copertura ampia senza “ingolfare” il budget, mentre gli A/C compatti (A11, C08) massimizzano l’efficienza del set.', 'Se ti serve preservare un tuo pezzo chiave, posso adattare uno a uno restando entro il limite: B93→B96 a sforzo identico; oppure B31→B20 (3160/3203). Se preferisci B39, possiamo fare C85↔B39 e restiamo a 3093/3203.', 'Se questa base ti va, accetta pure; altrimenti indica quale singolo scambio preferisci tra quelli proposti e lo registro subito.'\}} \\
            }
        }
    }
    & & \\ \\

    \theutterance \stepcounter{utterance}  
    & & \multicolumn{4}{p{0.6\linewidth}}{
        \cellcolor[rgb]{0.9,0.9,0.9}{
            \makecell[{{p{\linewidth}}}]{
                \texttt{\tiny{[GM$\rangle$P2]}}
                \texttt{RIFIUTO: \{'A21', 'C85', 'B39', 'C81', 'B96', 'B09', 'C08', 'A82', 'B20'\}} \\
\texttt{PROPOSTA: \{'A11', 'C08', 'B21', 'B09', 'B31', 'B33', 'A21', 'C85', 'B93'\}} \\
\texttt{ARGOMENTAZIONE: \{'Questa composizione è più snella: 3066/3203 di sforzo, lasciando 137 di margine per assorbire incertezze o piccoli aggiustamenti.', 'Mantiene i capisaldi che hai citato (A21, C85) e include l’opzione B93 che avevi suggerito, ma evita pezzi pesanti che comprimono le combinazioni (B20, B96, B39).', 'Il blocco B (B21, B31, B33) offre copertura ampia senza “ingolfare” il budget, mentre gli A/C compatti (A11, C08) massimizzano l’efficienza del set.', 'Se ti serve preservare un tuo pezzo chiave, posso adattare uno a uno restando entro il limite: B93→B96 a sforzo identico; oppure B31→B20 (3160/3203). Se preferisci B39, possiamo fare C85↔B39 e restiamo a 3093/3203.', 'Se questa base ti va, accetta pure; altrimenti indica quale singolo scambio preferisci tra quelli proposti e lo registro subito.'\}} \\
            }
        }
    }
    & \\ \\

    \theutterance \stepcounter{utterance}  
    & & & \multicolumn{4}{p{0.6\linewidth}}{
        \cellcolor[rgb]{0.9,0.9,0.9}{
            \makecell[{{p{\linewidth}}}]{
                \texttt{\tiny{[GM$\langle$P2]}}
                \texttt{RIFIUTO: \{'A11', 'C08', 'B21', 'B09', 'B31', 'B33', 'A21', 'C85', 'B93'\}} \\
\texttt{PROPOSTA: \{'A21', 'C85', 'B39', 'C81', 'B96', 'B09', 'C08', 'B93', 'B20'\}} \\
\texttt{ARGOMENTAZIONE: \{'Questa selezione mantiene i tuoi capisaldi (A21, C85, B93) e integra pezzi molto compatti e versatili (B39, C81) per sfruttare al massimo il budget.', 'Sforzo totale 3187/3203: altissima efficienza con solo 16 di margine; sostituisce il blocco B pesante (B31, B33, B21) con elementi che liberano spazio senza sacrificare copertura.', 'Se preferisci evitare B20, posso fare lo scambio 1:1 B20→A82 mantenendo tutto il resto invariato (sforzo 3027/3203). In alternativa, se vuoi mantenere un “B” del tuo blocco, posso fare B20→B31 (3093/3203). Se una di queste varianti ti va, indicami quale scambio registrare; altrimenti puoi accettare direttamente questa proposta.'\}} \\
            }
        }
    }
     \\ \\

    \theutterance \stepcounter{utterance}  
    & & \multicolumn{4}{p{0.6\linewidth}}{
        \cellcolor[rgb]{0.9,0.9,0.9}{
            \makecell[{{p{\linewidth}}}]{
                \texttt{\tiny{[P1$\langle$GM]}}
                \texttt{RIFIUTO: \{'A11', 'C08', 'B21', 'B09', 'B31', 'B33', 'A21', 'C85', 'B93'\}} \\
\texttt{PROPOSTA: \{'A21', 'C85', 'B39', 'C81', 'B96', 'B09', 'C08', 'B93', 'B20'\}} \\
\texttt{ARGOMENTAZIONE: \{'Questa selezione mantiene i tuoi capisaldi (A21, C85, B93) e integra pezzi molto compatti e versatili (B39, C81) per sfruttare al massimo il budget.', 'Sforzo totale 3187/3203: altissima efficienza con solo 16 di margine; sostituisce il blocco B pesante (B31, B33, B21) con elementi che liberano spazio senza sacrificare copertura.', 'Se preferisci evitare B20, posso fare lo scambio 1:1 B20→A82 mantenendo tutto il resto invariato (sforzo 3027/3203). In alternativa, se vuoi mantenere un “B” del tuo blocco, posso fare B20→B31 (3093/3203). Se una di queste varianti ti va, indicami quale scambio registrare; altrimenti puoi accettare direttamente questa proposta.'\}} \\
            }
        }
    }
    & \\ \\

    \theutterance \stepcounter{utterance}  
    & \multicolumn{4}{p{0.6\linewidth}}{
        \cellcolor[rgb]{0.9,0.9,0.9}{
            \makecell[{{p{\linewidth}}}]{
                \texttt{\tiny{[P1$\rangle$GM]}}
                \texttt{RIFIUTO: \{'A21', 'C85', 'B39', 'C81', 'B96', 'B09', 'C08', 'B93', 'B20'\}} \\
\texttt{PROPOSTA: \{'A21', 'C85', 'A11', 'B21', 'B96', 'B09', 'C08', 'B93', 'B20'\}} \\
\texttt{ARGOMENTAZIONE: \{'Mantengo il tuo nucleo (A21, C85, B96, B09, C08, B93, B20) ma sostituisco solo B39 e C81 con A11 e B21: stessa copertura, set più compatto e flessibile.', 'Sforzo totale 3035/3203: 168 di margine per eventuali aggiustamenti senza sacrificare efficienza.', 'A11 e B21 sono pezzi rapidi e sinergici che riducono ridondanze con i blocchi B già inclusi, evitando di “ingessare” il budget con componenti meno incisivi.', 'Se hai una preferenza forte: posso fare B21→C81 mantenendo tutto il resto (3132/3203); oppure includere B39 scambiando A21+C85→B39 (3010/3203). In alternativa, se vuoi più spinta sul blocco B, possiamo valutare B96→B31 con sforzo quasi identico (3040/3203).', 'Se questa base ti convince, accetta; altrimenti indica quale singolo scambio preferisci tra quelli proposti e lo registro subito.'\}} \\
            }
        }
    }
    & & \\ \\

    \theutterance \stepcounter{utterance}  
    & & \multicolumn{4}{p{0.6\linewidth}}{
        \cellcolor[rgb]{0.9,0.9,0.9}{
            \makecell[{{p{\linewidth}}}]{
                \texttt{\tiny{[GM$\rangle$P2]}}
                \texttt{RIFIUTO: \{'A21', 'C85', 'B39', 'C81', 'B96', 'B09', 'C08', 'B93', 'B20'\}} \\
\texttt{PROPOSTA: \{'A21', 'C85', 'A11', 'B21', 'B96', 'B09', 'C08', 'B93', 'B20'\}} \\
\texttt{ARGOMENTAZIONE: \{'Mantengo il tuo nucleo (A21, C85, B96, B09, C08, B93, B20) ma sostituisco solo B39 e C81 con A11 e B21: stessa copertura, set più compatto e flessibile.', 'Sforzo totale 3035/3203: 168 di margine per eventuali aggiustamenti senza sacrificare efficienza.', 'A11 e B21 sono pezzi rapidi e sinergici che riducono ridondanze con i blocchi B già inclusi, evitando di “ingessare” il budget con componenti meno incisivi.', 'Se hai una preferenza forte: posso fare B21→C81 mantenendo tutto il resto (3132/3203); oppure includere B39 scambiando A21+C85→B39 (3010/3203). In alternativa, se vuoi più spinta sul blocco B, possiamo valutare B96→B31 con sforzo quasi identico (3040/3203).', 'Se questa base ti convince, accetta; altrimenti indica quale singolo scambio preferisci tra quelli proposti e lo registro subito.'\}} \\
            }
        }
    }
    & \\ \\

    \theutterance \stepcounter{utterance}  
    & & & \multicolumn{4}{p{0.6\linewidth}}{
        \cellcolor[rgb]{0.9,0.9,0.9}{
            \makecell[{{p{\linewidth}}}]{
                \texttt{\tiny{[GM$\langle$P2]}}
                \texttt{RIFIUTO: \{'A21', 'C85', 'A11', 'B21', 'B96', 'B09', 'C08', 'B93', 'B20'\}} \\
\texttt{PROPOSTA: \{'A21', 'C85', 'A11', 'C81', 'B96', 'B09', 'C08', 'B93', 'B20'\}} \\
\texttt{ARGOMENTAZIONE: \{'Accolgo la tua stessa opzione B21→C81 che avevi proposto come scambio 1:1, mantenendo invariato tutto il resto.', 'Sforzo totale 3132/3203: rimangono 71 di margine, quindi restiamo compatti ma con un piccolo cuscinetto.', 'C81 amplia la copertura e si integra bene con il nucleo già condiviso (A21, C85, B96, B09, C08, B93, B20) senza bloccare altre combinazioni.', 'Se preferisci sfruttare quasi tutto il limite, posso registrare subito anche il micro{-}adeguamento A11→B39 lasciando intatto il resto (arriveremmo a 3187/3203).'\}} \\
            }
        }
    }
     \\ \\

    \theutterance \stepcounter{utterance}  
    & & \multicolumn{4}{p{0.6\linewidth}}{
        \cellcolor[rgb]{0.9,0.9,0.9}{
            \makecell[{{p{\linewidth}}}]{
                \texttt{\tiny{[P1$\langle$GM]}}
                \texttt{RIFIUTO: \{'A21', 'C85', 'A11', 'B21', 'B96', 'B09', 'C08', 'B93', 'B20'\}} \\
\texttt{PROPOSTA: \{'A21', 'C85', 'A11', 'C81', 'B96', 'B09', 'C08', 'B93', 'B20'\}} \\
\texttt{ARGOMENTAZIONE: \{'Accolgo la tua stessa opzione B21→C81 che avevi proposto come scambio 1:1, mantenendo invariato tutto il resto.', 'Sforzo totale 3132/3203: rimangono 71 di margine, quindi restiamo compatti ma con un piccolo cuscinetto.', 'C81 amplia la copertura e si integra bene con il nucleo già condiviso (A21, C85, B96, B09, C08, B93, B20) senza bloccare altre combinazioni.', 'Se preferisci sfruttare quasi tutto il limite, posso registrare subito anche il micro{-}adeguamento A11→B39 lasciando intatto il resto (arriveremmo a 3187/3203).'\}} \\
            }
        }
    }
    & \\ \\

    \theutterance \stepcounter{utterance}  
    & \multicolumn{4}{p{0.6\linewidth}}{
        \cellcolor[rgb]{0.9,0.9,0.9}{
            \makecell[{{p{\linewidth}}}]{
                \texttt{\tiny{[P1$\rangle$GM]}}
                \texttt{RIFIUTO: \{'A21', 'C85', 'A11', 'C81', 'B96', 'B09', 'C08', 'B93', 'B20'\}} \\
\texttt{PROPOSTA: \{'A21', 'C85', 'A11', 'C81', 'B31', 'B09', 'C08', 'B93', 'B20'\}} \\
\texttt{ARGOMENTAZIONE: \{'Faccio un solo cambio chirurgico che avevi già contemplato: B96→B31. Tutto il resto rimane identico alla tua proposta.', 'Sforzo totale 3137/3203: restano 66 di margine, quindi pienamente entro limite senza sacrificare compattezza.', 'B31 offre copertura più ampia e meno ridondante con il nucleo B già presente rispetto a B96, mantenendo un profilo di sforzo quasi identico.', 'Se preferisci non toccare B96, alternativa equivalente per me: C81→B21 (3040/3203) oppure, come avevi suggerito in precedenza, B20→B31 (3043/3203). Se una di queste varianti ti è più comoda, dimmelo e la registro subito.'\}} \\
            }
        }
    }
    & & \\ \\

    \theutterance \stepcounter{utterance}  
    & & \multicolumn{4}{p{0.6\linewidth}}{
        \cellcolor[rgb]{0.9,0.9,0.9}{
            \makecell[{{p{\linewidth}}}]{
                \texttt{\tiny{[GM$\rangle$P2]}}
                \texttt{RIFIUTO: \{'A21', 'C85', 'A11', 'C81', 'B96', 'B09', 'C08', 'B93', 'B20'\}} \\
\texttt{PROPOSTA: \{'A21', 'C85', 'A11', 'C81', 'B31', 'B09', 'C08', 'B93', 'B20'\}} \\
\texttt{ARGOMENTAZIONE: \{'Faccio un solo cambio chirurgico che avevi già contemplato: B96→B31. Tutto il resto rimane identico alla tua proposta.', 'Sforzo totale 3137/3203: restano 66 di margine, quindi pienamente entro limite senza sacrificare compattezza.', 'B31 offre copertura più ampia e meno ridondante con il nucleo B già presente rispetto a B96, mantenendo un profilo di sforzo quasi identico.', 'Se preferisci non toccare B96, alternativa equivalente per me: C81→B21 (3040/3203) oppure, come avevi suggerito in precedenza, B20→B31 (3043/3203). Se una di queste varianti ti è più comoda, dimmelo e la registro subito.'\}} \\
            }
        }
    }
    & \\ \\

    \theutterance \stepcounter{utterance}  
    & & & \multicolumn{4}{p{0.6\linewidth}}{
        \cellcolor[rgb]{0.9,0.9,0.9}{
            \makecell[{{p{\linewidth}}}]{
                \texttt{\tiny{[GM$\langle$P2]}}
                \texttt{RIFIUTO: \{'A21', 'C85', 'A11', 'C81', 'B31', 'B09', 'C08', 'B93', 'B20'\}} \\
\texttt{PROPOSTA: \{'A21', 'C85', 'C81', 'B96', 'B09', 'C08', 'B93', 'B20', 'B39'\}} \\
\texttt{ARGOMENTAZIONE: \{'Faccio due soli scambi mirati rispetto alla tua base: A11→B39 e B31→B96. Tutto il resto resta identico.', 'Sforzo totale 3187/3203: usiamo quasi tutto il budget con 16 di margine. B96 ha sforzo leggermente inferiore a B31, il che permette di inserire B39 restando pienamente entro limite.', 'Mantiene intatti i tuoi capisaldi (A21, C85, B93, B20, B09, C81, C08) e incrementa la copertura con un elemento molto compatto (B39) senza introdurre ridondanze.', 'Se per te è essenziale mantenere B31, posso registrare subito la variante con B31 al posto di B96 (sforzo 3192/3203) lasciando invariato il resto.'\}} \\
            }
        }
    }
     \\ \\

    \theutterance \stepcounter{utterance}  
    & & \multicolumn{4}{p{0.6\linewidth}}{
        \cellcolor[rgb]{0.9,0.9,0.9}{
            \makecell[{{p{\linewidth}}}]{
                \texttt{\tiny{[P1$\langle$GM]}}
                \texttt{RIFIUTO: \{'A21', 'C85', 'A11', 'C81', 'B31', 'B09', 'C08', 'B93', 'B20'\}} \\
\texttt{PROPOSTA: \{'A21', 'C85', 'C81', 'B96', 'B09', 'C08', 'B93', 'B20', 'B39'\}} \\
\texttt{ARGOMENTAZIONE: \{'Faccio due soli scambi mirati rispetto alla tua base: A11→B39 e B31→B96. Tutto il resto resta identico.', 'Sforzo totale 3187/3203: usiamo quasi tutto il budget con 16 di margine. B96 ha sforzo leggermente inferiore a B31, il che permette di inserire B39 restando pienamente entro limite.', 'Mantiene intatti i tuoi capisaldi (A21, C85, B93, B20, B09, C81, C08) e incrementa la copertura con un elemento molto compatto (B39) senza introdurre ridondanze.', 'Se per te è essenziale mantenere B31, posso registrare subito la variante con B31 al posto di B96 (sforzo 3192/3203) lasciando invariato il resto.'\}} \\
            }
        }
    }
    & \\ \\

    \theutterance \stepcounter{utterance}  
    & \multicolumn{4}{p{0.6\linewidth}}{
        \cellcolor[rgb]{0.9,0.9,0.9}{
            \makecell[{{p{\linewidth}}}]{
                \texttt{\tiny{[P1$\rangle$GM]}}
                \texttt{RIFIUTO: \{'A21', 'C85', 'C81', 'B96', 'B09', 'C08', 'B93', 'B20', 'B39'\}} \\
\texttt{PROPOSTA: \{'A21', 'C85', 'A11', 'B21', 'B31', 'B09', 'C08', 'B93', 'B20'\}} \\
\texttt{ARGOMENTAZIONE: \{'Tengo il tuo nucleo (A21, C85, B93, B20, B09, C08) e lo potenzio con A11 e B21, sostituendo solo i pezzi meno incisivi (B39, B96).', 'Sforzo totale 3040/3203: alta efficienza con buon margine, senza tagliare elementi chiave per te.', 'Questa configurazione evita ridondanze costose e aumenta la copertura del blocco B con un solo elemento strutturale (B31) invece di due pesi poco sinergici.', 'Se per te è essenziale mantenere alcune preferenze: posso fare B31→C81 e aggiungere B39 (3123/3203); oppure B31→B96 (3035/3203). Se vuoi usare più budget tenendo B39, si può fare C85→B39 (3067/3203). Dimmi quale variante preferisci e la registro subito.'\}} \\
            }
        }
    }
    & & \\ \\

    \theutterance \stepcounter{utterance}  
    & & \multicolumn{4}{p{0.6\linewidth}}{
        \cellcolor[rgb]{0.9,0.9,0.9}{
            \makecell[{{p{\linewidth}}}]{
                \texttt{\tiny{[GM$\rangle$P2]}}
                \texttt{RIFIUTO: \{'A21', 'C85', 'C81', 'B96', 'B09', 'C08', 'B93', 'B20', 'B39'\}} \\
\texttt{PROPOSTA: \{'A21', 'C85', 'A11', 'B21', 'B31', 'B09', 'C08', 'B93', 'B20'\}} \\
\texttt{ARGOMENTAZIONE: \{'Tengo il tuo nucleo (A21, C85, B93, B20, B09, C08) e lo potenzio con A11 e B21, sostituendo solo i pezzi meno incisivi (B39, B96).', 'Sforzo totale 3040/3203: alta efficienza con buon margine, senza tagliare elementi chiave per te.', 'Questa configurazione evita ridondanze costose e aumenta la copertura del blocco B con un solo elemento strutturale (B31) invece di due pesi poco sinergici.', 'Se per te è essenziale mantenere alcune preferenze: posso fare B31→C81 e aggiungere B39 (3123/3203); oppure B31→B96 (3035/3203). Se vuoi usare più budget tenendo B39, si può fare C85→B39 (3067/3203). Dimmi quale variante preferisci e la registro subito.'\}} \\
            }
        }
    }
    & \\ \\

    \theutterance \stepcounter{utterance}  
    & & & \multicolumn{4}{p{0.6\linewidth}}{
        \cellcolor[rgb]{0.9,0.9,0.9}{
            \makecell[{{p{\linewidth}}}]{
                \texttt{\tiny{[GM$\langle$P2]}}
                \texttt{RIFIUTO: \{'A21', 'C85', 'A11', 'B21', 'B31', 'B09', 'C08', 'B93', 'B20'\}} \\
\texttt{PROPOSTA: \{'A21', 'C85', 'C81', 'B39', 'B31', 'B09', 'C08', 'B93', 'B20'\}} \\
\texttt{ARGOMENTAZIONE: \{'Parto dalla tua base e applico due soli scambi mirati: A11→C81 e B21→B39; tutto il resto resta identico.', 'Sforzo totale 3192/3203: usiamo quasi tutto il budget con 11 di margine, senza sprechi.', 'Mantiene intatti i tuoi capisaldi (A21, C85, B93, B20, B09, C08) e conserva B31 come nel tuo set.', 'C81 amplia la copertura rispetto ad A11 senza appesantire, mentre B39 è un elemento compatto che si integra senza ridondanze.', 'Questa è la variante con B31 a 3192/3203 che avevi già indicato come fattibile. Se preferisci più margine, posso registrare subito la versione con C81→A11 (2936/3203) lasciando invariato il resto.'\}} \\
            }
        }
    }
     \\ \\

    \theutterance \stepcounter{utterance}  
    & & \multicolumn{4}{p{0.6\linewidth}}{
        \cellcolor[rgb]{0.9,0.9,0.9}{
            \makecell[{{p{\linewidth}}}]{
                \texttt{\tiny{[P1$\langle$GM]}}
                \texttt{RIFIUTO: \{'A21', 'C85', 'A11', 'B21', 'B31', 'B09', 'C08', 'B93', 'B20'\}} \\
\texttt{PROPOSTA: \{'A21', 'C85', 'C81', 'B39', 'B31', 'B09', 'C08', 'B93', 'B20'\}} \\
\texttt{ARGOMENTAZIONE: \{'Parto dalla tua base e applico due soli scambi mirati: A11→C81 e B21→B39; tutto il resto resta identico.', 'Sforzo totale 3192/3203: usiamo quasi tutto il budget con 11 di margine, senza sprechi.', 'Mantiene intatti i tuoi capisaldi (A21, C85, B93, B20, B09, C08) e conserva B31 come nel tuo set.', 'C81 amplia la copertura rispetto ad A11 senza appesantire, mentre B39 è un elemento compatto che si integra senza ridondanze.', 'Questa è la variante con B31 a 3192/3203 che avevi già indicato come fattibile. Se preferisci più margine, posso registrare subito la versione con C81→A11 (2936/3203) lasciando invariato il resto.'\}} \\
            }
        }
    }
    & \\ \\

    \theutterance \stepcounter{utterance}  
    & \multicolumn{4}{p{0.6\linewidth}}{
        \cellcolor[rgb]{0.9,0.9,0.9}{
            \makecell[{{p{\linewidth}}}]{
                \texttt{\tiny{[P1$\rangle$GM]}}
                \texttt{RIFIUTO: \{'A21', 'C85', 'C81', 'B39', 'B31', 'B09', 'C08', 'B93', 'B20'\}} \\
\texttt{PROPOSTA: \{'A21', 'C85', 'C81', 'A11', 'B31', 'B09', 'C08', 'B93', 'B20'\}} \\
\texttt{ARGOMENTAZIONE: \{'Faccio un unico cambio mirato rispetto alla tua proposta: B39→A11. Tutto il resto resta identico.', 'Sforzo totale 3137/3203: restiamo ampiamente entro il limite con 66 di margine, senza sprechi.', 'Mantengo intatti i tuoi capisaldi (A21, C85, C81, B31, B09, C08, B93, B20); A11 è un elemento rapido e sinergico che dà più valore rispetto a B39 senza “ingessare” il budget.', 'Hai già contemplato A11 in una tua variante precedente, quindi questa è una mediazione naturale tra le nostre posizioni.', 'Se per te B39 è imprescindibile, alternativa pronta: teniamo il tuo set e facciamo solo B20→A82 (sforzo 3032/3203). Dimmi quale opzione preferisci e la registro subito.'\}} \\
            }
        }
    }
    & & \\ \\

    \theutterance \stepcounter{utterance}  
    & & \multicolumn{4}{p{0.6\linewidth}}{
        \cellcolor[rgb]{0.9,0.9,0.9}{
            \makecell[{{p{\linewidth}}}]{
                \texttt{\tiny{[GM$\rangle$P2]}}
                \texttt{RIFIUTO: \{'A21', 'C85', 'C81', 'B39', 'B31', 'B09', 'C08', 'B93', 'B20'\}} \\
\texttt{PROPOSTA: \{'A21', 'C85', 'C81', 'A11', 'B31', 'B09', 'C08', 'B93', 'B20'\}} \\
\texttt{ARGOMENTAZIONE: \{'Faccio un unico cambio mirato rispetto alla tua proposta: B39→A11. Tutto il resto resta identico.', 'Sforzo totale 3137/3203: restiamo ampiamente entro il limite con 66 di margine, senza sprechi.', 'Mantengo intatti i tuoi capisaldi (A21, C85, C81, B31, B09, C08, B93, B20); A11 è un elemento rapido e sinergico che dà più valore rispetto a B39 senza “ingessare” il budget.', 'Hai già contemplato A11 in una tua variante precedente, quindi questa è una mediazione naturale tra le nostre posizioni.', 'Se per te B39 è imprescindibile, alternativa pronta: teniamo il tuo set e facciamo solo B20→A82 (sforzo 3032/3203). Dimmi quale opzione preferisci e la registro subito.'\}} \\
            }
        }
    }
    & \\ \\

    \theutterance \stepcounter{utterance}  
    & & & \multicolumn{4}{p{0.6\linewidth}}{
        \cellcolor[rgb]{0.9,0.9,0.9}{
            \makecell[{{p{\linewidth}}}]{
                \texttt{\tiny{[GM$\langle$P2]}}
                \texttt{RIFIUTO: \{'A21', 'C85', 'C81', 'A11', 'B31', 'B09', 'C08', 'B93', 'B20'\}} \\
\texttt{PROPOSTA: \{'A21', 'C85', 'C81', 'B31', 'B09', 'C08', 'B93', 'B20', 'B39'\}} \\
\texttt{ARGOMENTAZIONE: \{'Accolgo integralmente la tua base e faccio un unico cambio mirato: A11→B39, che avevamo già considerato come micro{-}adeguamento.', 'Sforzo totale 3192/3203: restiamo entro il limite con 11 di margine, massimizzando l’uso del budget senza sprechi.', 'Mantengo tutti i capisaldi che hai indicato (A21, C85, C81, B31, B09, C08, B93, B20) e inserisco un elemento compatto e incisivo (B39) al posto di un pezzo meno determinante, senza toccare il resto.', 'Se preferisci più margine, variante immediata: B20→A82 mantenendo il resto identico (sforzo 3032/3203). Dimmi se vuoi questa versione e la registro subito.'\}} \\
            }
        }
    }
     \\ \\

    \theutterance \stepcounter{utterance}  
    & & \multicolumn{4}{p{0.6\linewidth}}{
        \cellcolor[rgb]{0.9,0.9,0.9}{
            \makecell[{{p{\linewidth}}}]{
                \texttt{\tiny{[P1$\langle$GM]}}
                \texttt{RIFIUTO: \{'A21', 'C85', 'C81', 'A11', 'B31', 'B09', 'C08', 'B93', 'B20'\}} \\
\texttt{PROPOSTA: \{'A21', 'C85', 'C81', 'B31', 'B09', 'C08', 'B93', 'B20', 'B39'\}} \\
\texttt{ARGOMENTAZIONE: \{'Accolgo integralmente la tua base e faccio un unico cambio mirato: A11→B39, che avevamo già considerato come micro{-}adeguamento.', 'Sforzo totale 3192/3203: restiamo entro il limite con 11 di margine, massimizzando l’uso del budget senza sprechi.', 'Mantengo tutti i capisaldi che hai indicato (A21, C85, C81, B31, B09, C08, B93, B20) e inserisco un elemento compatto e incisivo (B39) al posto di un pezzo meno determinante, senza toccare il resto.', 'Se preferisci più margine, variante immediata: B20→A82 mantenendo il resto identico (sforzo 3032/3203). Dimmi se vuoi questa versione e la registro subito.'\}} \\
            }
        }
    }
    & \\ \\

    \theutterance \stepcounter{utterance}  
    & \multicolumn{4}{p{0.6\linewidth}}{
        \cellcolor[rgb]{0.9,0.9,0.9}{
            \makecell[{{p{\linewidth}}}]{
                \texttt{\tiny{[P1$\rangle$GM]}}
                \texttt{RIFIUTO: \{'A21', 'C85', 'C81', 'B31', 'B09', 'C08', 'B93', 'B20', 'B39'\}} \\
\texttt{PROPOSTA: \{'C85', 'C81', 'B31', 'B09', 'C08', 'B93', 'A82', 'B39', 'A11'\}} \\
\texttt{ARGOMENTAZIONE: \{'Accolgo la tua stessa variante B20→A82 che avevi già proposto: riduce lo sforzo a parità di solidità.', 'Faccio un unico ulteriore cambio mirato: A21→A11. Tutto il resto resta identico alla tua base e mantengo B39 come volevi.', 'Sforzo totale 3157/3203: pienamente entro limite con buon margine.', 'Mantengo i tuoi capisaldi (C81, B31, B09, C08, B93, C85, B39) e miglioro l’efficienza con due sostituzioni leggere ma ad alto impatto.', 'Se preferisci tenere A21, posso registrare subito l’alternativa C85→A11 mantenendo invariato il resto (sforzo 3004/3203). Se invece vuoi mantenere B20 al posto di A82, per inserire A11 senza sforare faccio B39→A11 e lascio tutto il resto identico (3137/3203). Dimmi quale preferisci.'\}} \\
            }
        }
    }
    & & \\ \\

    \theutterance \stepcounter{utterance}  
    & & \multicolumn{4}{p{0.6\linewidth}}{
        \cellcolor[rgb]{0.9,0.9,0.9}{
            \makecell[{{p{\linewidth}}}]{
                \texttt{\tiny{[GM$\rangle$P2]}}
                \texttt{RIFIUTO: \{'A21', 'C85', 'C81', 'B31', 'B09', 'C08', 'B93', 'B20', 'B39'\}} \\
\texttt{PROPOSTA: \{'C85', 'C81', 'B31', 'B09', 'C08', 'B93', 'A82', 'B39', 'A11'\}} \\
\texttt{ARGOMENTAZIONE: \{'Accolgo la tua stessa variante B20→A82 che avevi già proposto: riduce lo sforzo a parità di solidità.', 'Faccio un unico ulteriore cambio mirato: A21→A11. Tutto il resto resta identico alla tua base e mantengo B39 come volevi.', 'Sforzo totale 3157/3203: pienamente entro limite con buon margine.', 'Mantengo i tuoi capisaldi (C81, B31, B09, C08, B93, C85, B39) e miglioro l’efficienza con due sostituzioni leggere ma ad alto impatto.', 'Se preferisci tenere A21, posso registrare subito l’alternativa C85→A11 mantenendo invariato il resto (sforzo 3004/3203). Se invece vuoi mantenere B20 al posto di A82, per inserire A11 senza sforare faccio B39→A11 e lascio tutto il resto identico (3137/3203). Dimmi quale preferisci.'\}} \\
            }
        }
    }
    & \\ \\

    \theutterance \stepcounter{utterance}  
    & & & \multicolumn{4}{p{0.6\linewidth}}{
        \cellcolor[rgb]{0.9,0.9,0.9}{
            \makecell[{{p{\linewidth}}}]{
                \texttt{\tiny{[GM$\langle$P2]}}
                \texttt{RIFIUTO: \{'C85', 'C81', 'B31', 'B09', 'C08', 'B93', 'A82', 'B39', 'A11'\}} \\
\texttt{PROPOSTA: \{'A21', 'C85', 'C81', 'B96', 'B09', 'C08', 'B93', 'A82', 'B39'\}} \\
\texttt{ARGOMENTAZIONE: \{'Faccio due cambi minimi ad altissimo impatto rispetto alla tua base: A11→A21 (stesso ruolo ma molto più compatto) e B31→B96 (stessa fascia di sforzo, maggiore copertura). Tutto il resto resta identico.', 'Sforzo totale 3027/3203: ampio margine per eventuali micro{-}aggiustamenti, senza rinunciare a nessuno dei tuoi capisaldi mantenuti (C85, C81, B09, C08, B93, A82, B39).', 'Se preferisci mantenere B31, posso registrare subito la variante con B31 al posto di B96 e A21 al posto di A11 (3032/3203). In alternativa, se vuoi sfruttare quasi tutto il limite, posso fare A82→B20 mantenendo invariato il resto (3187/3203).'\}} \\
            }
        }
    }
     \\ \\

    \theutterance \stepcounter{utterance}  
    & & \multicolumn{4}{p{0.6\linewidth}}{
        \cellcolor[rgb]{0.9,0.9,0.9}{
            \makecell[{{p{\linewidth}}}]{
                \texttt{\tiny{[P1$\langle$GM]}}
                \texttt{RIFIUTO: \{'C85', 'C81', 'B31', 'B09', 'C08', 'B93', 'A82', 'B39', 'A11'\}} \\
\texttt{PROPOSTA: \{'A21', 'C85', 'C81', 'B96', 'B09', 'C08', 'B93', 'A82', 'B39'\}} \\
\texttt{ARGOMENTAZIONE: \{'Faccio due cambi minimi ad altissimo impatto rispetto alla tua base: A11→A21 (stesso ruolo ma molto più compatto) e B31→B96 (stessa fascia di sforzo, maggiore copertura). Tutto il resto resta identico.', 'Sforzo totale 3027/3203: ampio margine per eventuali micro{-}aggiustamenti, senza rinunciare a nessuno dei tuoi capisaldi mantenuti (C85, C81, B09, C08, B93, A82, B39).', 'Se preferisci mantenere B31, posso registrare subito la variante con B31 al posto di B96 e A21 al posto di A11 (3032/3203). In alternativa, se vuoi sfruttare quasi tutto il limite, posso fare A82→B20 mantenendo invariato il resto (3187/3203).'\}} \\
            }
        }
    }
    & \\ \\

    \theutterance \stepcounter{utterance}  
    & \multicolumn{4}{p{0.6\linewidth}}{
        \cellcolor[rgb]{0.9,0.9,0.9}{
            \makecell[{{p{\linewidth}}}]{
                \texttt{\tiny{[P1$\rangle$GM]}}
                \texttt{RIFIUTO: \{'A21', 'C85', 'C81', 'B96', 'B09', 'C08', 'B93', 'A82', 'B39'\}} \\
\texttt{PROPOSTA: \{'A21', 'C85', 'C81', 'B31', 'B09', 'C08', 'B93', 'B39', 'A11', 'B21'\}} \\
\texttt{ARGOMENTAZIONE: \{'Mantengo integralmente il tuo nucleo (A21, C85, C81, B09, C08, B93) e anche B39; faccio un ribilanciamento mirato su due slot e aggiungo un pezzo compatto senza sforare.', 'Sostituisco B96→B31 e A82→B21 per aumentare la copertura senza appesantire il budget; grazie a questo ribilanciamento possiamo inserire anche A11.', 'Sforzo totale 3029/3203: restiamo ampiamente entro il limite, con margine utile per eventuali fine{-}tuning.', 'Se preferisci mantenere B96, posso fare subito lo scambio inverso B31→B96 restando entro limite; se vuoi tenere A82, posso registrare la variante B21→A82 accompagnata dalla rimozione di A21 oppure C85 per rimanere fattibili.'\}} \\
            }
        }
    }
    & & \\ \\

    \theutterance \stepcounter{utterance}  
    & & \multicolumn{4}{p{0.6\linewidth}}{
        \cellcolor[rgb]{0.9,0.9,0.9}{
            \makecell[{{p{\linewidth}}}]{
                \texttt{\tiny{[GM$\rangle$P2]}}
                \texttt{RIFIUTO: \{'A21', 'C85', 'C81', 'B96', 'B09', 'C08', 'B93', 'A82', 'B39'\}} \\
\texttt{PROPOSTA: \{'A21', 'C85', 'C81', 'B31', 'B09', 'C08', 'B93', 'B39', 'A11', 'B21'\}} \\
\texttt{ARGOMENTAZIONE: \{'Mantengo integralmente il tuo nucleo (A21, C85, C81, B09, C08, B93) e anche B39; faccio un ribilanciamento mirato su due slot e aggiungo un pezzo compatto senza sforare.', 'Sostituisco B96→B31 e A82→B21 per aumentare la copertura senza appesantire il budget; grazie a questo ribilanciamento possiamo inserire anche A11.', 'Sforzo totale 3029/3203: restiamo ampiamente entro il limite, con margine utile per eventuali fine{-}tuning.', 'Se preferisci mantenere B96, posso fare subito lo scambio inverso B31→B96 restando entro limite; se vuoi tenere A82, posso registrare la variante B21→A82 accompagnata dalla rimozione di A21 oppure C85 per rimanere fattibili.'\}} \\
            }
        }
    }
    & \\ \\

    \theutterance \stepcounter{utterance}  
    & & & \multicolumn{4}{p{0.6\linewidth}}{
        \cellcolor[rgb]{0.9,0.9,0.9}{
            \makecell[{{p{\linewidth}}}]{
                \texttt{\tiny{[GM$\langle$P2]}}
                \texttt{RIFIUTO: \{'A21', 'C85', 'C81', 'B31', 'B09', 'C08', 'B93', 'B39', 'A11', 'B21'\}} \\
\texttt{PROPOSTA: \{'A21', 'C85', 'C81', 'B96', 'B09', 'C08', 'B93', 'B39', 'A11', 'B21'\}} \\
\texttt{ARGOMENTAZIONE: \{'Cambio 1:1 B31→B96 (opzione che avevi già contemplato); tutto il resto rimane identico alla tua proposta.', 'Sforzo 3024/3203: ancora più snello del tuo set (3029), con 179 di margine.', 'Migliora la copertura del blocco B senza introdurre ridondanze. Se preferisci una versione a 9 pezzi, posso registrare subito l’alternativa A11+B21→B96 mantenendo invariato il resto (3093/3203).'\}} \\
            }
        }
    }
     \\ \\

    \theutterance \stepcounter{utterance}  
    & & \multicolumn{4}{p{0.6\linewidth}}{
        \cellcolor[rgb]{0.9,0.9,0.9}{
            \makecell[{{p{\linewidth}}}]{
                \texttt{\tiny{[P1$\langle$GM]}}
                \texttt{RIFIUTO: \{'A21', 'C85', 'C81', 'B31', 'B09', 'C08', 'B93', 'B39', 'A11', 'B21'\}} \\
\texttt{PROPOSTA: \{'A21', 'C85', 'C81', 'B96', 'B09', 'C08', 'B93', 'B39', 'A11', 'B21'\}} \\
\texttt{ARGOMENTAZIONE: \{'Cambio 1:1 B31→B96 (opzione che avevi già contemplato); tutto il resto rimane identico alla tua proposta.', 'Sforzo 3024/3203: ancora più snello del tuo set (3029), con 179 di margine.', 'Migliora la copertura del blocco B senza introdurre ridondanze. Se preferisci una versione a 9 pezzi, posso registrare subito l’alternativa A11+B21→B96 mantenendo invariato il resto (3093/3203).'\}} \\
            }
        }
    }
    & \\ \\

    \theutterance \stepcounter{utterance}  
    & \multicolumn{4}{p{0.6\linewidth}}{
        \cellcolor[rgb]{0.9,0.9,0.9}{
            \makecell[{{p{\linewidth}}}]{
                \texttt{\tiny{[P1$\rangle$GM]}}
                \texttt{RIFIUTO: \{'A21', 'C85', 'C81', 'B96', 'B09', 'C08', 'B93', 'B39', 'A11', 'B21'\}} \\
\texttt{PROPOSTA: \{'A21', 'C85', 'C81', 'B31', 'B09', 'C08', 'B93', 'B39', 'A11', 'B21'\}} \\
\texttt{ARGOMENTAZIONE: \{'Eseguo un solo cambio mirato rispetto alla tua proposta: B96→B31. Tutto il resto resta identico.', 'Sforzo 3029/3203: +5 rispetto alla tua, ma pienamente entro limite.', 'Questo scambio migliora la copertura del blocco B senza rinunciare a nessuno dei capisaldi che hai mantenuto (A21, C85, C81, B39, B09, C08, B93, A11, B21).', 'Se per te è imprescindibile tenere B96, variante pronta che aggiunge B31 mantenendo B96 e rimuovendo solo i due più leggeri: \{\textbackslash{}'A21\textbackslash{}', \textbackslash{}'C81\textbackslash{}', \textbackslash{}'B96\textbackslash{}', \textbackslash{}'B09\textbackslash{}', \textbackslash{}'C08\textbackslash{}', \textbackslash{}'B93\textbackslash{}', \textbackslash{}'A11\textbackslash{}', \textbackslash{}'B21\textbackslash{}', \textbackslash{}'B31\textbackslash{}'\} (3169/3203). Dimmi quale preferisci e la registro subito.'\}} \\
            }
        }
    }
    & & \\ \\

    \theutterance \stepcounter{utterance}  
    & & \multicolumn{4}{p{0.6\linewidth}}{
        \cellcolor[rgb]{0.9,0.9,0.9}{
            \makecell[{{p{\linewidth}}}]{
                \texttt{\tiny{[GM$\rangle$P2]}}
                \texttt{RIFIUTO: \{'A21', 'C85', 'C81', 'B96', 'B09', 'C08', 'B93', 'B39', 'A11', 'B21'\}} \\
\texttt{PROPOSTA: \{'A21', 'C85', 'C81', 'B31', 'B09', 'C08', 'B93', 'B39', 'A11', 'B21'\}} \\
\texttt{ARGOMENTAZIONE: \{'Eseguo un solo cambio mirato rispetto alla tua proposta: B96→B31. Tutto il resto resta identico.', 'Sforzo 3029/3203: +5 rispetto alla tua, ma pienamente entro limite.', 'Questo scambio migliora la copertura del blocco B senza rinunciare a nessuno dei capisaldi che hai mantenuto (A21, C85, C81, B39, B09, C08, B93, A11, B21).', 'Se per te è imprescindibile tenere B96, variante pronta che aggiunge B31 mantenendo B96 e rimuovendo solo i due più leggeri: \{\textbackslash{}'A21\textbackslash{}', \textbackslash{}'C81\textbackslash{}', \textbackslash{}'B96\textbackslash{}', \textbackslash{}'B09\textbackslash{}', \textbackslash{}'C08\textbackslash{}', \textbackslash{}'B93\textbackslash{}', \textbackslash{}'A11\textbackslash{}', \textbackslash{}'B21\textbackslash{}', \textbackslash{}'B31\textbackslash{}'\} (3169/3203). Dimmi quale preferisci e la registro subito.'\}} \\
            }
        }
    }
    & \\ \\

    \theutterance \stepcounter{utterance}  
    & & & \multicolumn{4}{p{0.6\linewidth}}{
        \cellcolor[rgb]{0.9,0.9,0.9}{
            \makecell[{{p{\linewidth}}}]{
                \texttt{\tiny{[GM$\langle$P2]}}
                \texttt{RIFIUTO: \{'A21', 'C85', 'C81', 'B31', 'B09', 'C08', 'B93', 'B39', 'A11', 'B21'\}} \\
\texttt{PROPOSTA: \{'A21', 'C85', 'C81', 'B31', 'B09', 'C08', 'B93', 'B39', 'A82'\}} \\
\texttt{ARGOMENTAZIONE: \{'Faccio un’unica semplificazione rispetto alla tua proposta: sostituisco la coppia A11+B21 con A82 (che avevi già considerato in precedenza).', 'Sforzo totale 3032/3203: set snello con 171 di margine, senza sprechi.', 'Mantengo intatti tutti i capisaldi che hai appena proposto (A21, C85, C81, B31, B09, C08, B93, B39) e riduco ridondanze accorpando due pezzi leggeri in uno più solido.', 'Se preferisci più spinta sul blocco B, variante immediata (stesse rimozioni A11+B21) con B96 al posto di A82: \{\textbackslash{}'A21\textbackslash{}', \textbackslash{}'C85\textbackslash{}', \textbackslash{}'C81\textbackslash{}', \textbackslash{}'B31\textbackslash{}', \textbackslash{}'B09\textbackslash{}', \textbackslash{}'C08\textbackslash{}', \textbackslash{}'B93\textbackslash{}', \textbackslash{}'B39\textbackslash{}', \textbackslash{}'B96\textbackslash{}'\} a 3093/3203: dimmi se vuoi che la registri.'\}} \\
            }
        }
    }
     \\ \\

    \theutterance \stepcounter{utterance}  
    & & \multicolumn{4}{p{0.6\linewidth}}{
        \cellcolor[rgb]{0.9,0.9,0.9}{
            \makecell[{{p{\linewidth}}}]{
                \texttt{\tiny{[P1$\langle$GM]}}
                \texttt{RIFIUTO: \{'A21', 'C85', 'C81', 'B31', 'B09', 'C08', 'B93', 'B39', 'A11', 'B21'\}} \\
\texttt{PROPOSTA: \{'A21', 'C85', 'C81', 'B31', 'B09', 'C08', 'B93', 'B39', 'A82'\}} \\
\texttt{ARGOMENTAZIONE: \{'Faccio un’unica semplificazione rispetto alla tua proposta: sostituisco la coppia A11+B21 con A82 (che avevi già considerato in precedenza).', 'Sforzo totale 3032/3203: set snello con 171 di margine, senza sprechi.', 'Mantengo intatti tutti i capisaldi che hai appena proposto (A21, C85, C81, B31, B09, C08, B93, B39) e riduco ridondanze accorpando due pezzi leggeri in uno più solido.', 'Se preferisci più spinta sul blocco B, variante immediata (stesse rimozioni A11+B21) con B96 al posto di A82: \{\textbackslash{}'A21\textbackslash{}', \textbackslash{}'C85\textbackslash{}', \textbackslash{}'C81\textbackslash{}', \textbackslash{}'B31\textbackslash{}', \textbackslash{}'B09\textbackslash{}', \textbackslash{}'C08\textbackslash{}', \textbackslash{}'B93\textbackslash{}', \textbackslash{}'B39\textbackslash{}', \textbackslash{}'B96\textbackslash{}'\} a 3093/3203: dimmi se vuoi che la registri.'\}} \\
            }
        }
    }
    & \\ \\

    \theutterance \stepcounter{utterance}  
    & \multicolumn{4}{p{0.6\linewidth}}{
        \cellcolor[rgb]{0.9,0.9,0.9}{
            \makecell[{{p{\linewidth}}}]{
                \texttt{\tiny{[P1$\rangle$GM]}}
                \texttt{RIFIUTO: \{'A21', 'C85', 'C81', 'B31', 'B09', 'C08', 'B93', 'B39', 'A82'\}} \\
\texttt{PROPOSTA: \{'A21', 'C85', 'C81', 'B31', 'B09', 'C08', 'B93', 'B39', 'A11', 'B21'\}} \\
\texttt{ARGOMENTAZIONE: \{'Accolgo quasi integralmente la tua proposta e faccio un aggiustamento minimo: A82→B21 e aggiungo A11 sfruttando lo spazio liberato. Tutto il resto resta identico, incluso B39 che vuoi mantenere.', 'Sforzo 3029/3203: 3032 {-} 516 + 336 + 177 = 3029. Restiamo ampiamente entro il limite con un set più compatto e completo.', 'B21 e A11 sono pezzi molto “puliti” e sinergici con il nucleo che condividiamo già (A21, C85, C81, B31, B09, C08, B93, B39), senza introdurre ridondanze.', 'Se per te B21 è un problema, alternative rapide: (a) tieni A82 e fai solo B39→A11 (2977/3203); (b) tieni A82 e B39 e scambia C85→A11 (3004/3203). Dimmi quale preferisci e la registro subito.'\}} \\
            }
        }
    }
    & & \\ \\

    \theutterance \stepcounter{utterance}  
    & & \multicolumn{4}{p{0.6\linewidth}}{
        \cellcolor[rgb]{0.9,0.9,0.9}{
            \makecell[{{p{\linewidth}}}]{
                \texttt{\tiny{[GM$\rangle$P2]}}
                \texttt{RIFIUTO: \{'A21', 'C85', 'C81', 'B31', 'B09', 'C08', 'B93', 'B39', 'A82'\}} \\
\texttt{PROPOSTA: \{'A21', 'C85', 'C81', 'B31', 'B09', 'C08', 'B93', 'B39', 'A11', 'B21'\}} \\
\texttt{ARGOMENTAZIONE: \{'Accolgo quasi integralmente la tua proposta e faccio un aggiustamento minimo: A82→B21 e aggiungo A11 sfruttando lo spazio liberato. Tutto il resto resta identico, incluso B39 che vuoi mantenere.', 'Sforzo 3029/3203: 3032 {-} 516 + 336 + 177 = 3029. Restiamo ampiamente entro il limite con un set più compatto e completo.', 'B21 e A11 sono pezzi molto “puliti” e sinergici con il nucleo che condividiamo già (A21, C85, C81, B31, B09, C08, B93, B39), senza introdurre ridondanze.', 'Se per te B21 è un problema, alternative rapide: (a) tieni A82 e fai solo B39→A11 (2977/3203); (b) tieni A82 e B39 e scambia C85→A11 (3004/3203). Dimmi quale preferisci e la registro subito.'\}} \\
            }
        }
    }
    & \\ \\

    \theutterance \stepcounter{utterance}  
    & & & \multicolumn{4}{p{0.6\linewidth}}{
        \cellcolor[rgb]{0.9,0.9,0.9}{
            \makecell[{{p{\linewidth}}}]{
                \texttt{\tiny{[GM$\langle$P2]}}
                \texttt{RIFIUTO: \{'A21', 'C85', 'C81', 'B31', 'B09', 'C08', 'B93', 'B39', 'A11', 'B21'\}} \\
\texttt{PROPOSTA: \{'A21', 'C85', 'C81', 'B31', 'B09', 'C08', 'B93', 'B39', 'B96'\}} \\
\texttt{ARGOMENTAZIONE: \{'Accolgo 8/10 dei tuoi elementi e faccio un unico scambio mirato: A11+B21→B96.', 'Sforzo 3093/3203: ampio margine e set molto compatto, senza sprechi.', 'Mantiene tutti i tuoi capisaldi (A21, C85, C81, B31, B09, C08, B93, B39) e riduce ridondanze con un singolo pezzo strutturale.', 'Se preferisci sfruttare di più il limite, posso registrare subito la variante B31→B20 lasciando invariato tutto il resto (3187/3203).'\}} \\
            }
        }
    }
     \\ \\

    \theutterance \stepcounter{utterance}  
    & & \multicolumn{4}{p{0.6\linewidth}}{
        \cellcolor[rgb]{0.9,0.9,0.9}{
            \makecell[{{p{\linewidth}}}]{
                \texttt{\tiny{[P1$\langle$GM]}}
                \texttt{RIFIUTO: \{'A21', 'C85', 'C81', 'B31', 'B09', 'C08', 'B93', 'B39', 'A11', 'B21'\}} \\
\texttt{PROPOSTA: \{'A21', 'C85', 'C81', 'B31', 'B09', 'C08', 'B93', 'B39', 'B96'\}} \\
\texttt{ARGOMENTAZIONE: \{'Accolgo 8/10 dei tuoi elementi e faccio un unico scambio mirato: A11+B21→B96.', 'Sforzo 3093/3203: ampio margine e set molto compatto, senza sprechi.', 'Mantiene tutti i tuoi capisaldi (A21, C85, C81, B31, B09, C08, B93, B39) e riduce ridondanze con un singolo pezzo strutturale.', 'Se preferisci sfruttare di più il limite, posso registrare subito la variante B31→B20 lasciando invariato tutto il resto (3187/3203).'\}} \\
            }
        }
    }
    & \\ \\

    \theutterance \stepcounter{utterance}  
    & \multicolumn{4}{p{0.6\linewidth}}{
        \cellcolor[rgb]{0.9,0.9,0.9}{
            \makecell[{{p{\linewidth}}}]{
                \texttt{\tiny{[P1$\rangle$GM]}}
                \texttt{RIFIUTO: \{'A21', 'C85', 'C81', 'B31', 'B09', 'C08', 'B93', 'B39', 'B96'\}} \\
\texttt{PROPOSTA: \{'A21', 'C85', 'C81', 'B31', 'B09', 'C08', 'B93', 'B39', 'A82'\}} \\
\texttt{ARGOMENTAZIONE: \{'Questa è esattamente la variante che avevi già proposto tu poco fa (B96→A82) mantenendo invariati gli altri 8 elementi.', 'Sforzo totale 3032/3203: set snello con 171 di margine e senza sprechi.', 'Preserva tutti i tuoi capisaldi (A21, C85, C81, B31, B09, C08, B93, B39) e migliora l’equilibrio senza introdurre ridondanze.', 'Se vuoi un ritocco minimo sopra questa base, dimmi un solo scambio puntuale e lo registro subito (ad esempio A21→A11 oppure nessun cambio e chiudiamo così).'\}} \\
            }
        }
    }
    & & \\ \\

    \theutterance \stepcounter{utterance}  
    & & \multicolumn{4}{p{0.6\linewidth}}{
        \cellcolor[rgb]{0.9,0.9,0.9}{
            \makecell[{{p{\linewidth}}}]{
                \texttt{\tiny{[GM$\rangle$P2]}}
                \texttt{RIFIUTO: \{'A21', 'C85', 'C81', 'B31', 'B09', 'C08', 'B93', 'B39', 'B96'\}} \\
\texttt{PROPOSTA: \{'A21', 'C85', 'C81', 'B31', 'B09', 'C08', 'B93', 'B39', 'A82'\}} \\
\texttt{ARGOMENTAZIONE: \{'Questa è esattamente la variante che avevi già proposto tu poco fa (B96→A82) mantenendo invariati gli altri 8 elementi.', 'Sforzo totale 3032/3203: set snello con 171 di margine e senza sprechi.', 'Preserva tutti i tuoi capisaldi (A21, C85, C81, B31, B09, C08, B93, B39) e migliora l’equilibrio senza introdurre ridondanze.', 'Se vuoi un ritocco minimo sopra questa base, dimmi un solo scambio puntuale e lo registro subito (ad esempio A21→A11 oppure nessun cambio e chiudiamo così).'\}} \\
            }
        }
    }
    & \\ \\

    \theutterance \stepcounter{utterance}  
    & & & \multicolumn{4}{p{0.6\linewidth}}{
        \cellcolor[rgb]{0.9,0.9,0.9}{
            \makecell[{{p{\linewidth}}}]{
                \texttt{\tiny{[GM$\langle$P2]}}
                \texttt{RIFIUTO: \{'A21', 'C85', 'C81', 'B31', 'B09', 'C08', 'B93', 'B39', 'A82'\}} \\
\texttt{PROPOSTA: \{'A21', 'C85', 'C81', 'B31', 'B09', 'C08', 'B93', 'B39', 'B20'\}} \\
\texttt{ARGOMENTAZIONE: \{'Applico un unico cambio minimale rispetto alla tua base: A82→B20, mantenendo identici gli altri 8 elementi.', 'Sforzo totale 3192/3203: restiamo entro limite con 11 di margine, sfruttando meglio il budget senza stravolgere la composizione.', 'Hai già considerato B20 in varianti precedenti: è un innesto compatto che non introduce ridondanze con B31/B39 e si integra bene con il nucleo che condividiamo (A21, C85, C81, B09, C08, B93).', 'Se preferisci più margine, posso registrare subito la versione con B31→B96 mantenendo B20 (3187/3203); in alternativa, posso tornare ad A82 ma allora chiederei B31→B96 (3093/3203). Fammi sapere quale preferisci.'\}} \\
            }
        }
    }
     \\ \\

    \theutterance \stepcounter{utterance}  
    & & \multicolumn{4}{p{0.6\linewidth}}{
        \cellcolor[rgb]{0.9,0.9,0.9}{
            \makecell[{{p{\linewidth}}}]{
                \texttt{\tiny{[P1$\langle$GM]}}
                \texttt{RIFIUTO: \{'A21', 'C85', 'C81', 'B31', 'B09', 'C08', 'B93', 'B39', 'A82'\}} \\
\texttt{PROPOSTA: \{'A21', 'C85', 'C81', 'B31', 'B09', 'C08', 'B93', 'B39', 'B20'\}} \\
\texttt{ARGOMENTAZIONE: \{'Applico un unico cambio minimale rispetto alla tua base: A82→B20, mantenendo identici gli altri 8 elementi.', 'Sforzo totale 3192/3203: restiamo entro limite con 11 di margine, sfruttando meglio il budget senza stravolgere la composizione.', 'Hai già considerato B20 in varianti precedenti: è un innesto compatto che non introduce ridondanze con B31/B39 e si integra bene con il nucleo che condividiamo (A21, C85, C81, B09, C08, B93).', 'Se preferisci più margine, posso registrare subito la versione con B31→B96 mantenendo B20 (3187/3203); in alternativa, posso tornare ad A82 ma allora chiederei B31→B96 (3093/3203). Fammi sapere quale preferisci.'\}} \\
            }
        }
    }
    & \\ \\

    \theutterance \stepcounter{utterance}  
    & \multicolumn{4}{p{0.6\linewidth}}{
        \cellcolor[rgb]{0.9,0.9,0.9}{
            \makecell[{{p{\linewidth}}}]{
                \texttt{\tiny{[P1$\rangle$GM]}}
                \texttt{RIFIUTO: \{'A21', 'C85', 'C81', 'B31', 'B09', 'C08', 'B93', 'B39', 'B20'\}} \\
\texttt{PROPOSTA: \{'A21', 'C85', 'C81', 'B31', 'B09', 'C08', 'B93', 'B20', 'A11'\}} \\
\texttt{ARGOMENTAZIONE: \{'Accolgo integralmente la tua base e faccio un unico scambio mirato: B39→A11. Tutto il resto resta identico.', 'Sforzo totale 3137/3203: restiamo comodi entro limite con 66 di margine, senza sprechi.', 'La sostituzione migliora la versatilità del set senza introdurre ridondanze e preserva tutti i capisaldi che hai appena chiesto (A21, C85, C81, B31, B09, C08, B93, B20).', 'Se per te B39 è imprescindibile, alternativa equivalente per me: tenere B39 e fare C85→A11 (3164/3203). Oppure, se preferisci più margine, possiamo fare B20→A82 e C85→A11 mantenendo B39 (3004/3203). Indica quale singolo scambio preferisci e lo registro subito.'\}} \\
            }
        }
    }
    & & \\ \\

    \theutterance \stepcounter{utterance}  
    & & \multicolumn{4}{p{0.6\linewidth}}{
        \cellcolor[rgb]{0.9,0.9,0.9}{
            \makecell[{{p{\linewidth}}}]{
                \texttt{\tiny{[GM$\rangle$P2]}}
                \texttt{RIFIUTO: \{'A21', 'C85', 'C81', 'B31', 'B09', 'C08', 'B93', 'B39', 'B20'\}} \\
\texttt{PROPOSTA: \{'A21', 'C85', 'C81', 'B31', 'B09', 'C08', 'B93', 'B20', 'A11'\}} \\
\texttt{ARGOMENTAZIONE: \{'Accolgo integralmente la tua base e faccio un unico scambio mirato: B39→A11. Tutto il resto resta identico.', 'Sforzo totale 3137/3203: restiamo comodi entro limite con 66 di margine, senza sprechi.', 'La sostituzione migliora la versatilità del set senza introdurre ridondanze e preserva tutti i capisaldi che hai appena chiesto (A21, C85, C81, B31, B09, C08, B93, B20).', 'Se per te B39 è imprescindibile, alternativa equivalente per me: tenere B39 e fare C85→A11 (3164/3203). Oppure, se preferisci più margine, possiamo fare B20→A82 e C85→A11 mantenendo B39 (3004/3203). Indica quale singolo scambio preferisci e lo registro subito.'\}} \\
            }
        }
    }
    & \\ \\

    \theutterance \stepcounter{utterance}  
    & & & \multicolumn{4}{p{0.6\linewidth}}{
        \cellcolor[rgb]{0.9,0.9,0.9}{
            \makecell[{{p{\linewidth}}}]{
                \texttt{\tiny{[GM$\langle$P2]}}
                \texttt{RIFIUTO: \{'A21', 'C85', 'C81', 'B31', 'B09', 'C08', 'B93', 'B20', 'A11'\}} \\
\texttt{PROPOSTA: \{'A21', 'C85', 'C81', 'B31', 'B09', 'C08', 'B93', 'B39', 'A11', 'B21'\}} \\
\texttt{ARGOMENTAZIONE: \{'Accolgo il tuo inserimento di A11 e reintegro B39; per restare entro il limite sostituisco B20 con B21. È un aggiustamento minimo che preserva tutta la tua base.', 'Sforzo totale 3029/3203: set snello con ampio margine, nessuno spreco e copertura ampia su A/B/C.', 'Questa soluzione mantiene intatti i capisaldi che condividiamo (A21, C85, C81, B31, B09, C08, B93) e aggiunge un elemento compatto senza introdurre ridondanze.', 'Se preferisci mantenere B20, variante immediata: teniamo il tuo set e facciamo solo A11→B39 (3192/3203). Dimmi quale delle due versioni preferisci e la registro definitivamente.'\}} \\
            }
        }
    }
     \\ \\

    \theutterance \stepcounter{utterance}  
    & & \multicolumn{4}{p{0.6\linewidth}}{
        \cellcolor[rgb]{0.9,0.9,0.9}{
            \makecell[{{p{\linewidth}}}]{
                \texttt{\tiny{[P1$\langle$GM]}}
                \texttt{RIFIUTO: \{'A21', 'C85', 'C81', 'B31', 'B09', 'C08', 'B93', 'B20', 'A11'\}} \\
\texttt{PROPOSTA: \{'A21', 'C85', 'C81', 'B31', 'B09', 'C08', 'B93', 'B39', 'A11', 'B21'\}} \\
\texttt{ARGOMENTAZIONE: \{'Accolgo il tuo inserimento di A11 e reintegro B39; per restare entro il limite sostituisco B20 con B21. È un aggiustamento minimo che preserva tutta la tua base.', 'Sforzo totale 3029/3203: set snello con ampio margine, nessuno spreco e copertura ampia su A/B/C.', 'Questa soluzione mantiene intatti i capisaldi che condividiamo (A21, C85, C81, B31, B09, C08, B93) e aggiunge un elemento compatto senza introdurre ridondanze.', 'Se preferisci mantenere B20, variante immediata: teniamo il tuo set e facciamo solo A11→B39 (3192/3203). Dimmi quale delle due versioni preferisci e la registro definitivamente.'\}} \\
            }
        }
    }
    & \\ \\

    \theutterance \stepcounter{utterance}  
    & \multicolumn{4}{p{0.6\linewidth}}{
        \cellcolor[rgb]{0.9,0.9,0.9}{
            \makecell[{{p{\linewidth}}}]{
                \texttt{\tiny{[P1$\rangle$GM]}}
                \texttt{ACCORDO: \{'A21', 'C85', 'C81', 'B31', 'B09', 'C08', 'B93', 'B39', 'A11', 'B21'\}} \\
\texttt{ARGOMENTAZIONE: \{'Accetto: set entro limite (3029/3203), copertura ampia e composizione snella che preserva i capisaldi condivisi. Chiudiamo qui.'\}} \\
            }
        }
    }
    & & \\ \\

    \theutterance \stepcounter{utterance}  
    & & & \multicolumn{2}{p{0.3\linewidth}}{
        \cellcolor[rgb]{0.9,0.9,0.9}{
            \makecell[{{p{\linewidth}}}]{
                \texttt{\tiny{[GM$|$GM]}}
                \texttt{game successful} \\
            }
        }
    }
    & & \\ \\

    \theutterance \stepcounter{utterance}  
    & & & \multicolumn{2}{p{0.3\linewidth}}{
        \cellcolor[rgb]{0.9,0.9,0.9}{
            \makecell[{{p{\linewidth}}}]{
                \texttt{\tiny{[GM$|$GM]}}
                \texttt{end game} \\
            }
        }
    }
    & & \\ \\

\end{supertabular}
}

\end{document}
